\subsection{小 Batch 改变竞争区间}
\label{sec:7-smallbatch}

T5 在四个 E1 数据集、batch 8/32 上比较调度 JHQ、CAGRA fp32／int8 和 IVF-RaBitQ。五轮同步主机到主机测量分别使用 1,000/992 查询。1,808 个测点包含 JHQ 单 block 对照，不含 IVF-PQ／暴力搜索。\Cref{fig:7-smallbatch,tab:7-smallbatch} 比较所得区间。

\suppressfloats[t]
\par\medskip\noindent\begin{minipage}{\linewidth}
\captionsetup{hypcap=false}
\centering
\includegraphics[width=\linewidth]{t4-t5/smallbatch.pdf}
\captionof{figure}[小 Batch 竞争区间]{T5 的 JHQ／最快对手 QPS 比，采用保留合格实测端点的插值。大于一有利于 JHQ。严格阈值见 \Cref{tab:7-smallbatch}。}
\label{fig:7-smallbatch}
\end{minipage}\par\medskip

\par\medskip\noindent\begin{minipage}{\linewidth}
\captionsetup{hypcap=false}
\centering
\captionof{table}[小 Batch 匹配召回比]{T5 QPS 比：保留合格端点的插值／严格实测阈值。各估计器独立选择最快对手；满足阈值的更高召回端点可入选，无需外推。}
\label{tab:7-smallbatch}
\footnotesize\setlength{\tabcolsep}{4pt}
\begin{tabular}{@{}lrrrr@{}}
\toprule
数据集 & $b$ & $R=0.90$ & $R=0.95$ & $R=0.97$ \\
\midrule
Vogue & 8 & 1.155 / 1.091 & 1.073 / 0.980 & 0.893 / 0.854 \\
Vogue & 32 & 1.069 / 0.997 & 0.956 / 0.895 & 0.865 / 0.839 \\
arXiv & 8 & 1.479 / 1.454 & 0.859 / 0.851 & 0.697 / 0.698 \\
arXiv & 32 & 1.166 / 1.148 & 0.667 / 0.661 & 0.622 / 0.582 \\
OpenAI-1536 & 8 & 1.316 / 1.131 & 0.851 / 0.825 & 0.641 / 0.620 \\
OpenAI-1536 & 32 & 1.071 / 0.915 & 0.770 / 0.752 & 0.937 / 0.914 \\
OpenAI-3072 & 8 & 1.032 / 0.953 & 0.798 / 0.798 & 0.721 / 0.702 \\
OpenAI-3072 & 32 & 0.866 / 0.734 & 0.608 / 0.771 & 1.008 / 1.351 \\
\bottomrule
\end{tabular}

\end{minipage}\par\medskip

batch 8／召回率 0.97 时，CAGRA fp32 在两种估计器下均赢得全部四个数据集，JHQ 比为 0.641--0.893。batch 32 时比为 0.622--1.008；实测阈值下 OpenAI-3072 有利于 JHQ。Vogue 的 batch-8／0.95 领先需要插值。网格分辨率会改变接近持平的结果。

纳入合格端点时，召回率 0.90 比值为 batch 8 的 1.032--1.479、32 的 0.866--1.166；严格实测下 JHQ 分别赢三／一个。省略端点会夸大收益。调度在 0.95 处相对单 block 提升 $1.7$--$3.3\times$，但不保证跨系统获胜。

空调用中位数为 0.012--0.072\,ms，仍计入总时间。CAGRA 使用 Python，JHQ／RaBitQ 使用 C++；预热不同（\cref{sec:a-t4-t5}）。
